\section{C++ 与汇编互操作}

\subsection{本章要解决的问题}

上一章讲清楚了 System V AMD64 ABI：参数放在哪里、返回值放在哪里、哪些寄存器要保存、栈如何对齐。那些规则看起来像表格，但真正的价值在工程边界上才会显现出来：一个 C++ 文件如何调用手写汇编？手写汇编如何再调用 C++ helper？链接器如何把两个目标文件接到一起？内联汇编为什么经常比你想象中危险？为什么同样一条 \instruction{call}，在目标文件里还没有最终地址，链接之后才变成真实跳转？

本章把“能读汇编”推进到“能写可维护、可测试、可审计的汇编边界”。你需要掌握的不只是语法，而是一套工程流程：

\begin{lstlisting}
source declaration:
    C++ header says which symbol exists and what ABI contract it has

assembly implementation:
    .S file exports that symbol and obeys ABI

object file:
    assembler emits symbols, sections, relocations, debug/unwind hints

linking:
    linker resolves references and patches call/data addresses

runtime:
    CPU only executes bytes; correctness depends on the binary contract
\end{lstlisting}

从高性能计算和 AI 算子开发角度，本章非常关键。真实算子库经常出现这样的边界：

\begin{itemize}
  \item C++ 调度层选择某个 kernel，再调用汇编或 intrinsic 实现。
  \item 汇编 kernel 接收指针、维度、步长、常量、尾部长度等参数。
  \item kernel 内部尽量不调用其他函数；一旦调用，就必须恢复 ABI 约束。
  \item profiler、debugger、sanitizer、异常展开和符号工具都要能看懂这段代码。
\end{itemize}

因此，本章的核心不是“写几条汇编指令”，而是建立下面这条链路：

\begin{lstlisting}
C++ type/signature
    -> ABI-level parameter and return-value locations
    -> assembly symbol and prologue/epilogue
    -> object-file symbols and relocations
    -> linked executable
    -> debugger/profiler observable behavior
    -> tests and performance evidence
\end{lstlisting}

\begin{keyidea}
C++ 与汇编互操作的本质是二进制契约工程。C++ 编译器不会理解你手写汇编的意图；汇编器也不会检查你的 C++ 类型语义。两边能正确合作，完全依赖符号名、ABI、目标文件格式、链接器和测试共同维持边界。
\end{keyidea}

\subsection{先建立三层边界}

很多初学者写手写汇编失败，不是因为不会写 \instruction{add}，而是混淆了三层不同的规则。

\subsubsection{第一层：C++ 语言层}

C++ 源代码里看到的是函数声明、类型、重载、命名空间、引用、对象生命周期、异常、访问控制和模板实例化。例如：

\begin{lstlisting}[language=C++]
extern "C" long asm_sum8(long a, long b, long c, long d,
                         long e, long f, long g, long h);
\end{lstlisting}

这个声明告诉 C++ 编译器：

\begin{itemize}
  \item 存在一个可调用函数。
  \item 函数名在链接层叫 \code{asm_sum8}，不要使用 C++ name mangling。
  \item 参数和返回值都是 \code{long}。
  \item 调用时按当前平台默认调用约定生成代码。
\end{itemize}

但这个声明没有告诉编译器函数体在哪里，也没有证明汇编实现真的返回 \code{long}。如果你的汇编函数把结果放进了 \register{rcx} 而不是 \register{rax}，C++ 编译器不会替你发现。

\subsubsection{第二层：ABI 层}

上一章的 ABI 表格在这里变成实际契约。对于上面的 \code{asm_sum8}，System V AMD64 下入口状态应当是：

\begin{lstlisting}
on entry to asm_sum8:
    rdi       = a
    rsi       = b
    rdx       = c
    rcx       = d
    r8        = e
    r9        = f
    [rsp + 0] = return address
    [rsp + 8] = g
    [rsp +16] = h

on return:
    rax = result
\end{lstlisting}

如果函数修改了 \register{rbx}、\register{rbp}、\register{r12} 到 \register{r15}，它必须在返回前恢复。若它调用其他函数，调用点前还要满足栈对齐约定。

\subsubsection{第三层：目标文件和链接层}

汇编文件最终不是直接和 C++ 源码连接，而是先变成目标文件。目标文件里包含：

\begin{itemize}
  \item section，例如 \code{.text}、\code{.rodata}、\code{.data}。
  \item symbol，例如 \code{asm_sum8}、\code{cpp_helper}。
  \item relocation，例如“这里有一个 \instruction{call}，目标符号稍后由链接器填”。
  \item debug/unwind metadata，例如调试器和异常展开需要的信息。
\end{itemize}

可以用下面工具观察：

\begin{lstlisting}[language=bash]
nm -C file.o
readelf -Ws file.o
readelf -r file.o
objdump -drwC -Mintel file.o
\end{lstlisting}

\begin{mentalmodel}
源代码层问“这个函数叫什么、类型是什么”；ABI 层问“参数和返回值在哪些寄存器或栈槽”；目标文件层问“哪个符号定义在哪里，哪个地址需要链接器以后修补”。这三层都正确，跨语言调用才正确。
\end{mentalmodel}

\subsection{一个最小可运行工程}

先做最小例子。目录结构如下：

\begin{lstlisting}
labs/ch14_interop/minimal/
    CMakeLists.txt
    main.cpp
    asm_add.S
\end{lstlisting}

\subsubsection{C++ 调用者}

\begin{lstlisting}[language=C++]
// main.cpp
#include <cassert>
#include <cstdint>
#include <cstdio>

extern "C" std::int64_t asm_add2(std::int64_t a, std::int64_t b);

int main()
{
    const std::int64_t x = 40;
    const std::int64_t y = 2;
    const std::int64_t got = asm_add2(x, y);

    assert(got == 42);
    std::printf("asm_add2(%ld, %ld) = %ld\n", x, y, got);
}
\end{lstlisting}

这里的 \code{extern "C"} 不是说这个函数由 C 写成，而是说这个声明使用 C 链接名。没有它，C++ 可能把符号名改写成带类型信息的 mangled name。

\subsubsection{汇编实现}

\begin{lstlisting}
// asm_add.S
.intel_syntax noprefix
.text

.globl asm_add2
.type asm_add2, @function
asm_add2:
    lea rax, [rdi + rsi]
    ret
.size asm_add2, .-asm_add2

.section .note.GNU-stack,"",@progbits
\end{lstlisting}

逐行解释：

\begin{center}
\begin{tabular}{p{0.30\linewidth}p{0.55\linewidth}}
\toprule
代码 & 含义 \\
\midrule
\code{.intel_syntax noprefix} & 使用 Intel 语法，寄存器不写 \code{\%} 前缀。 \\
\code{.text} & 后续内容进入代码段。 \\
\code{.globl asm_add2} & 导出符号，让其他目标文件可以引用。 \\
\code{.type asm_add2, @function} & 标记符号类型为函数，便于工具识别。 \\
\code{asm_add2:} & 函数入口标签，也是符号地址。 \\
\code{lea rax, [rdi + rsi]} & 计算第 1、2 个参数之和，放入返回寄存器。 \\
\code{ret} & 返回到调用者。 \\
\code{.size asm_add2, .-asm_add2} & 记录函数大小，从入口到当前位置。 \\
\code{.note.GNU-stack} & 告诉链接器不需要可执行栈。 \\
\bottomrule
\end{tabular}
\end{center}

\begin{warningbox}
\code{.section .note.GNU-stack,"",@progbits} 不会生成运行时代码，但在 Linux ELF 工程中应当保留。省略它可能导致链接器给出 executable stack 警告，或者在某些环境里产生不必要的安全风险。
\end{warningbox}

\subsubsection{构建脚本}

\begin{lstlisting}
# CMakeLists.txt
cmake_minimum_required(VERSION 3.20)
project(ch14_minimal LANGUAGES CXX ASM)

add_executable(ch14_minimal
    main.cpp
    asm_add.S
)

target_compile_features(ch14_minimal PRIVATE cxx_std_20)
target_compile_options(ch14_minimal PRIVATE
    $<$<COMPILE_LANGUAGE:CXX>:-Wall -Wextra -Wpedantic -O2 -g>
)
\end{lstlisting}

构建：

\begin{lstlisting}[language=bash]
cmake -S . -B build -DCMAKE_BUILD_TYPE=RelWithDebInfo
cmake --build build -j
./build/ch14_minimal
\end{lstlisting}

审计：

\begin{lstlisting}[language=bash]
nm -C build/ch14_minimal | rg 'asm_add2|main'
objdump -drwC -Mintel build/ch14_minimal > ch14_minimal.objdump.txt
rg -n '<asm_add2>|call.*asm_add2' ch14_minimal.objdump.txt
\end{lstlisting}

\subsubsection{为什么文件后缀用 \code{.S}}

GNU/Clang 工具链里常见约定：

\begin{center}
\begin{tabular}{p{0.16\linewidth}p{0.68\linewidth}}
\toprule
后缀 & 含义 \\
\midrule
\code{.s} & 汇编源文件，通常不经过 C 预处理器。 \\
\code{.S} & 汇编源文件，先经过 C 预处理器，再交给汇编器。 \\
\bottomrule
\end{tabular}
\end{center}

使用 \code{.S} 后，可以写：

\begin{lstlisting}
#if defined(__linux__) && defined(__x86_64__)
.section .note.GNU-stack,"",@progbits
#endif
\end{lstlisting}

也可以在以后用宏生成不同 ISA 版本。但宏会降低可读性。本书训练阶段优先写直白的汇编，只有重复结构明显时才引入宏。

\subsection{把函数签名当成二进制接口}

互操作时，最危险的习惯是“先写汇编，能跑就行”。正确顺序应该是：

\begin{enumerate}
  \item 写 C++ 声明，明确类型和链接名。
  \item 按 ABI 推导参数位置、返回位置、保存责任。
  \item 写汇编函数框架。
  \item 写测试，覆盖普通值、边界值、随机值。
  \item 用 \code{objdump}、\code{nm}、\code{readelf} 审计二进制。
  \item 再考虑性能。
\end{enumerate}

例如：

\begin{lstlisting}[language=C++]
extern "C" std::int64_t asm_sum8(std::int64_t a0, std::int64_t a1,
                                 std::int64_t a2, std::int64_t a3,
                                 std::int64_t a4, std::int64_t a5,
                                 std::int64_t a6, std::int64_t a7);
\end{lstlisting}

ABI 推导表：

\begin{center}
\begin{tabular}{p{0.18\linewidth}p{0.24\linewidth}p{0.42\linewidth}}
\toprule
源代码参数 & 入口位置 & 说明 \\
\midrule
\code{a0} & \register{rdi} & 第 1 个整数类参数。 \\
\code{a1} & \register{rsi} & 第 2 个整数类参数。 \\
\code{a2} & \register{rdx} & 第 3 个整数类参数。 \\
\code{a3} & \register{rcx} & 第 4 个整数类参数。 \\
\code{a4} & \register{r8} & 第 5 个整数类参数。 \\
\code{a5} & \register{r9} & 第 6 个整数类参数。 \\
\code{a6} & \code{[rsp + 8]} & 第 7 个参数，返回地址之后。 \\
\code{a7} & \code{[rsp + 16]} & 第 8 个参数。 \\
return & \register{rax} & 64 位整数返回。 \\
\bottomrule
\end{tabular}
\end{center}

汇编：

\begin{lstlisting}
.intel_syntax noprefix
.text

.globl asm_sum8
.type asm_sum8, @function
asm_sum8:
    lea rax, [rdi + rsi]
    add rax, rdx
    add rax, rcx
    add rax, r8
    add rax, r9
    add rax, qword ptr [rsp + 8]
    add rax, qword ptr [rsp + 16]
    ret
.size asm_sum8, .-asm_sum8

.section .note.GNU-stack,"",@progbits
\end{lstlisting}

注意这里没有函数序言。它是一个 leaf function，不调用其他函数，不需要局部栈空间，不修改 callee-saved 寄存器。这样的函数可以非常短。但这个短小建立在 ABI 推导正确之上，不是随便省略。

\begin{deepdive}
如果函数入口后执行了 \instruction{push} \register{rbp}，那么第 7 个参数就不再位于当前 \register{rsp} 的 \code{[rsp + 8]}。入口原始布局被改变了：
\begin{lstlisting}
on entry:
    [rsp + 0]  return address
    [rsp + 8]  a6
    [rsp +16]  a7

after push rbp:
    [rsp + 0]  saved rbp
    [rsp + 8]  return address
    [rsp +16]  a6
    [rsp +24]  a7
\end{lstlisting}
所以读写栈参数前必须先确定当前 \register{rsp} 相对入口移动了多少。很多手写汇编 bug 就来自这里。
\end{deepdive}

\subsection{测试不是附属品}

手写汇编很容易出现“几个样例能跑，但边界条件错”的情况。测试要直接围绕二进制契约设计。

\subsubsection{确定性测试}

\begin{lstlisting}[language=C++]
#include <array>
#include <cassert>
#include <cstdint>
#include <limits>

extern "C" std::int64_t asm_sum8(std::int64_t, std::int64_t,
                                 std::int64_t, std::int64_t,
                                 std::int64_t, std::int64_t,
                                 std::int64_t, std::int64_t);

static std::int64_t ref_sum8(const std::array<std::int64_t, 8>& values)
{
    std::int64_t sum = 0;
    for (const std::int64_t value : values) {
        sum += value;
    }
    return sum;
}

int main()
{
    const std::array<std::array<std::int64_t, 8>, 4> cases {{
        {{0, 0, 0, 0, 0, 0, 0, 0}},
        {{1, 2, 3, 4, 5, 6, 7, 8}},
        {{-1, -2, -3, -4, -5, -6, -7, -8}},
        {{100, -50, 7, -8, 9, -10, 11, -12}},
    }};

    for (const auto& c : cases) {
        const std::int64_t got = asm_sum8(c[0], c[1], c[2], c[3],
                                          c[4], c[5], c[6], c[7]);
        assert(got == ref_sum8(c));
    }
}
\end{lstlisting}

这个测试至少覆盖了：

\begin{itemize}
  \item 寄存器参数。
  \item 栈参数 \code{a6} 和 \code{a7}。
  \item 正数、负数、零。
  \item 多次调用。
\end{itemize}

\subsubsection{随机测试}

\begin{lstlisting}[language=C++]
#include <array>
#include <cassert>
#include <cstdint>
#include <random>

extern "C" std::int64_t asm_sum8(std::int64_t, std::int64_t,
                                 std::int64_t, std::int64_t,
                                 std::int64_t, std::int64_t,
                                 std::int64_t, std::int64_t);

static std::int64_t ref_sum8(const std::array<std::int64_t, 8>& values)
{
    std::int64_t sum = 0;
    for (const std::int64_t value : values) {
        sum += value;
    }
    return sum;
}

int main()
{
    std::mt19937_64 rng(0xC001D00DULL);
    std::uniform_int_distribution<std::int64_t> dist(-1000000, 1000000);

    for (int iter = 0; iter < 100000; ++iter) {
        std::array<std::int64_t, 8> values {};
        for (std::int64_t& value : values) {
            value = dist(rng);
        }

        const std::int64_t got = asm_sum8(values[0], values[1],
                                          values[2], values[3],
                                          values[4], values[5],
                                          values[6], values[7]);
        assert(got == ref_sum8(values));
    }
}
\end{lstlisting}

\begin{warningbox}
如果 C++ reference 函数使用 signed overflow，而测试数据可能溢出，那么 reference 本身可能有未定义行为。训练早期可以限制输入范围；后续如果要测试完整 64 位 wraparound，应使用 \code{std::uint64_t} 或明确使用编译器支持的溢出内建函数。
\end{warningbox}

\subsection{C++ name mangling 与 \code{extern "C"}}

C++ 支持重载、命名空间、类成员函数、模板实例化。为了让链接器区分这些函数，编译器会把源代码函数名编码成更复杂的链接符号，这叫 \term{name mangling}{名称修饰}。

\begin{lstlisting}[language=C++]
int add(int a, int b)
{
    return a + b;
}

double add(double a, double b)
{
    return a + b;
}
\end{lstlisting}

编译并查看符号：

\begin{lstlisting}[language=bash]
clang++ -std=c++20 -O2 -c overload.cpp -o overload.o
nm overload.o
nm -C overload.o
\end{lstlisting}

可能看到类似：

\begin{lstlisting}
0000000000000000 T _Z3addii
0000000000000010 T _Z3adddd
\end{lstlisting}

\code{nm -C} 会 demangle：

\begin{lstlisting}
0000000000000000 T add(int, int)
0000000000000010 T add(double, double)
\end{lstlisting}

如果写：

\begin{lstlisting}[language=C++]
extern "C" int asm_add_int(int a, int b);
\end{lstlisting}

链接符号就是：

\begin{lstlisting}
asm_add_int
\end{lstlisting}

这让手写汇编更简单。汇编文件只需要导出同名符号：

\begin{lstlisting}
.globl asm_add_int
.type asm_add_int, @function
asm_add_int:
    lea eax, [edi + esi]
    ret
.size asm_add_int, .-asm_add_int
\end{lstlisting}

\begin{deepdive}
\code{extern "C"} 只改变语言链接规则和符号名，不会把 C++ 类型系统变成 C，也不会禁用 System V ABI。比如 \code{extern "C" double f(double)} 仍然用 \register{xmm0} 传参与返回。它也不能用于重载同名函数，因为 C 链接名无法表达重载集合。
\end{deepdive}

\subsection{目标文件里的符号和重定位}

很多人第一次看目标文件反汇编，会困惑为什么 \instruction{call} 后面是 0。

看 C++：

\begin{lstlisting}[language=C++]
extern "C" long asm_add2(long, long);

extern "C" long call_add2()
{
    return asm_add2(10, 32);
}
\end{lstlisting}

编译目标文件：

\begin{lstlisting}[language=bash]
clang++ -std=c++20 -O2 -c caller.cpp -o caller.o
objdump -drwC -Mintel caller.o
readelf -r caller.o
\end{lstlisting}

你可能看到：

\begin{lstlisting}
0000000000000000 <call_add2>:
   0: bf 0a 00 00 00        mov    edi,0xa
   5: be 20 00 00 00        mov    esi,0x20
   a: e8 00 00 00 00        call   f <call_add2+0xf>
        b: R_X86_64_PLT32   asm_add2-0x4
   f: c3                    ret
\end{lstlisting}

这里的 \code{e8 00 00 00 00} 不是最终机器码。x86 的 near \instruction{call} 编码通常是：

\begin{lstlisting}
e8 rel32
\end{lstlisting}

\code{rel32} 是相对下一条指令地址的有符号 32 位位移。目标文件阶段还不知道 \code{asm_add2} 最终地址，所以汇编器放一个临时值，并生成 relocation 记录：

\begin{lstlisting}
at offset 0xb:
    patch 4 bytes
    relocation type = R_X86_64_PLT32
    symbol = asm_add2
    addend = -4
\end{lstlisting}

链接后，假设：

\begin{lstlisting}
call instruction address = 0x40112a
next rip after call      = 0x40112f
asm_add2 address         = 0x401150
\end{lstlisting}

那么：

\begin{lstlisting}
rel32 = target - next_rip
      = 0x401150 - 0x40112f
      = 0x21
\end{lstlisting}

最终编码会接近：

\begin{lstlisting}
e8 21 00 00 00
\end{lstlisting}

\begin{keyidea}
目标文件里的反汇编不是最终地址事实，而是“机器码字节 + 符号 + 重定位请求”的组合。读未链接目标文件时，必须同时看 \code{objdump -dr} 和 \code{readelf -r}。
\end{keyidea}

\subsection{汇编函数调用 C++ 函数}

现在反过来：手写汇编调用 C++ helper。

\subsubsection{C++ helper}

\begin{lstlisting}[language=C++]
#include <cstdint>

extern "C" std::int64_t cpp_square(std::int64_t x)
{
    return x * x;
}

extern "C" std::int64_t asm_square_plus(std::int64_t x);
\end{lstlisting}

目标：\code{asm_square_plus(x)} 调用 \code{cpp_square(x)}，再返回 \code{x*x + x}。

错误版本很容易写成：

\begin{lstlisting}
.intel_syntax noprefix
.text

.globl asm_square_plus_broken
.type asm_square_plus_broken, @function
asm_square_plus_broken:
    mov rbx, rdi
    call cpp_square
    add rax, rbx
    ret
.size asm_square_plus_broken, .-asm_square_plus_broken
\end{lstlisting}

这段有两个风险：

\begin{itemize}
  \item 修改了 callee-saved \register{rbx}，但没有恢复。
  \item 进入函数时 \register{rsp} 通常是 \code{16n + 8}；直接 \instruction{call} 时调用点前栈没有 16 字节对齐。
\end{itemize}

\subsubsection{修复版本}

\begin{lstlisting}
.intel_syntax noprefix
.text

.globl asm_square_plus
.type asm_square_plus, @function
asm_square_plus:
    push rbx
    mov  rbx, rdi

    call cpp_square

    add  rax, rbx
    pop  rbx
    ret
.size asm_square_plus, .-asm_square_plus

.section .note.GNU-stack,"",@progbits
\end{lstlisting}

为什么一个 \instruction{push} \register{rbx} 同时解决了两个问题？

入口时：

\begin{lstlisting}
rsp = 16n + 8
\end{lstlisting}

执行 \instruction{push} \register{rbx} 后：

\begin{lstlisting}
rsp = 16n
\end{lstlisting}

这时执行 \instruction{call} \code{cpp_square}，调用点前 \register{rsp} 满足 16 字节对齐。同时 \register{rbx} 原值被保存到栈上，返回前 \instruction{pop} \register{rbx} 恢复。

栈变化：

\begin{lstlisting}
on entry to asm_square_plus:
    [rsp + 0]  return address to C++ caller

after push rbx:
    [rsp + 0]  saved rbx
    [rsp + 8]  return address to C++ caller

inside cpp_square after call:
    [rsp + 0]  return address to asm_square_plus
    [rsp + 8]  saved rbx
    [rsp +16]  return address to C++ caller
\end{lstlisting}

\begin{warningbox}
不要把“push 一个寄存器刚好对齐”当成万能模板。如果你还要分配局部栈空间，必须重新计算总共改变了多少 \register{rsp}。原则是：每个 \instruction{call} 执行前，调用者的 \register{rsp} 应为 16 字节对齐。
\end{warningbox}

\subsection{leaf 和 non-leaf 汇编模板}

手写汇编可以先从两个模板开始。

\subsubsection{leaf function 模板}

leaf function 不调用其他函数。它可以非常简洁：

\begin{lstlisting}
.intel_syntax noprefix
.text

.globl function_name
.type function_name, @function
function_name:
    # use caller-saved registers freely:
    # rax, rcx, rdx, rsi, rdi, r8-r11, xmm0-xmm15
    # preserve callee-saved registers if you touch them

    ret
.size function_name, .-function_name

.section .note.GNU-stack,"",@progbits
\end{lstlisting}

对于短小整数函数，常见策略是只用 caller-saved 寄存器，避免保存恢复开销。

\subsubsection{non-leaf function 模板}

non-leaf function 会调用其他函数，必须更谨慎：

\begin{lstlisting}
.intel_syntax noprefix
.text

.globl function_name
.type function_name, @function
function_name:
    push rbp
    mov  rbp, rsp
    push rbx
    sub  rsp, 8

    # now rsp is 16-byte aligned before calls
    # save live caller-saved values yourself before each call if needed

    call some_helper

    add  rsp, 8
    pop  rbx
    pop  rbp
    ret
.size function_name, .-function_name

.section .note.GNU-stack,"",@progbits
\end{lstlisting}

为什么有 \code{sub rsp, 8}？计算一下：

\begin{lstlisting}
entry:
    rsp = 16n + 8

push rbp:
    rsp = 16n

mov rbp, rsp:
    no change

push rbx:
    rsp = 16n - 8

sub rsp, 8:
    rsp = 16n - 16
\end{lstlisting}

所以在 \instruction{call} \code{some_helper} 前，\register{rsp} 是 16 字节对齐。

\begin{deepdive}
真实编译器会根据局部变量大小、保存寄存器数量、是否省略帧指针、是否需要栈保护、是否有异常展开信息等因素生成不同序言。模板不是背诵答案，而是帮助你建立栈对齐和保存责任的算术模型。
\end{deepdive}

\subsection{C++ 对象、引用和指针在边界上的形状}

汇编只看见寄存器、内存和地址。C++ 的引用、指针、小结构体、大结构体在 ABI 层有不同形状。

\subsubsection{引用通常就是地址}

C++：

\begin{lstlisting}[language=C++]
extern "C" void asm_inc_ref(std::int64_t& x);
\end{lstlisting}

虽然源代码写的是引用，但汇编入口看到的通常是指针：

\begin{lstlisting}
on entry:
    rdi = address of x
\end{lstlisting}

汇编实现：

\begin{lstlisting}
.globl asm_inc_ref
.type asm_inc_ref, @function
asm_inc_ref:
    inc qword ptr [rdi]
    ret
.size asm_inc_ref, .-asm_inc_ref
\end{lstlisting}

这不意味着 C++ 引用“语言上等于指针”。语言层引用有别名、生命周期和绑定规则；ABI 层只是经常用地址传递它。

\subsubsection{小结构体可能走寄存器}

\begin{lstlisting}[language=C++]
struct Pair64 {
    std::int64_t lo;
    std::int64_t hi;
};

extern "C" Pair64 asm_make_pair(std::int64_t x, std::int64_t y);
\end{lstlisting}

两个 64 位整数大小的结构体，在常见 System V AMD64 规则下可用 \register{rax} 和 \register{rdx} 返回：

\begin{lstlisting}
.globl asm_make_pair
.type asm_make_pair, @function
asm_make_pair:
    mov rax, rdi
    mov rdx, rsi
    ret
.size asm_make_pair, .-asm_make_pair
\end{lstlisting}

\subsubsection{大结构体通常通过隐藏指针返回}

\begin{lstlisting}[language=C++]
struct Triple64 {
    std::int64_t a;
    std::int64_t b;
    std::int64_t c;
};

extern "C" Triple64 asm_make_triple(std::int64_t x,
                                    std::int64_t y,
                                    std::int64_t z);
\end{lstlisting}

常见入口状态：

\begin{lstlisting}
rdi = hidden pointer to return object
rsi = x
rdx = y
rcx = z
\end{lstlisting}

实现：

\begin{lstlisting}
.globl asm_make_triple
.type asm_make_triple, @function
asm_make_triple:
    mov rax, rdi
    mov qword ptr [rdi + 0], rsi
    mov qword ptr [rdi + 8], rdx
    mov qword ptr [rdi +16], rcx
    ret
.size asm_make_triple, .-asm_make_triple
\end{lstlisting}

注意显式参数 \code{x} 不在 \register{rdi}，因为 \register{rdi} 被隐藏返回对象指针占用了。这是读跨语言汇编时很常见的坑。

\subsection{全局数据和 RIP-relative addressing}

在 x86-64 位置无关代码中，访问静态数据常用 RIP-relative addressing。

\begin{lstlisting}
.intel_syntax noprefix
.text

.globl asm_get_message
.type asm_get_message, @function
asm_get_message:
    lea rax, [rip + message]
    ret
.size asm_get_message, .-asm_get_message

.section .rodata
message:
    .asciz "hello from assembly"

.section .note.GNU-stack,"",@progbits
\end{lstlisting}

C++ 声明：

\begin{lstlisting}[language=C++]
extern "C" const char* asm_get_message();
\end{lstlisting}

\instruction{lea} 不是读取字符串内容，而是计算字符串地址：

\begin{lstlisting}
rax = address of message
\end{lstlisting}

反汇编目标文件时，仍然可能看到 relocation，因为 \code{message} 的最终地址要到链接阶段才知道：

\begin{lstlisting}
lea rax, [rip + 0x0]
    relocation: R_X86_64_PC32 .rodata-0x4
\end{lstlisting}

\begin{deepdive}
RIP-relative addressing 的基准是下一条指令的地址，不是当前指令开头地址。这个规则和 \instruction{call} \code{rel32} 一样，都是 x86-64 位置无关代码的重要基础。后面讲动态库、PLT/GOT 和 PIC 时会继续展开。
\end{deepdive}

\subsection{内联汇编为什么危险}

很多人一学汇编就想在 C++ 里写 inline asm。工程上要反过来：除非你非常清楚编译器优化器和寄存器分配器会怎样看它，否则优先使用普通 C++、compiler builtin、intrinsics 或单独的 \code{.S} 文件。

GNU extended asm 的基本形状：

\begin{lstlisting}[language=C++]
asm volatile (
    "assembly template"
    : output operands
    : input operands
    : clobbers
);
\end{lstlisting}

它不是“把几行汇编插入 C++”这么简单。它是在告诉优化器：

\begin{itemize}
  \item 哪些值被这个 asm 读。
  \item 哪些值被这个 asm 写。
  \item 哪些寄存器、flags、内存状态被破坏。
  \item 这个 asm 是否可以删除、移动或合并。
\end{itemize}

如果这些信息写错，编译器可能生成看似合理但实际错误的代码。

\subsubsection{错误示例：没有声明输出}

\begin{lstlisting}[language=C++]
std::int64_t broken_add(std::int64_t a, std::int64_t b)
{
    std::int64_t result = a;
    asm("add %1, %0" : : "r"(result), "r"(b));
    return result;
}
\end{lstlisting}

这段的问题是：模板里修改了 \code{\%0} 对应的寄存器，但约束里没有告诉编译器 \code{result} 是输出。优化器可以认为 \code{result} 没有改变。

修复：

\begin{lstlisting}[language=C++]
std::int64_t fixed_add(std::int64_t a, std::int64_t b)
{
    std::int64_t result = a;
    asm("add %1, %0"
        : "+r"(result)
        : "r"(b)
        : "cc");
    return result;
}
\end{lstlisting}

\code{"+r"(result)} 表示这个操作数既输入又输出，并且放在某个通用寄存器中。\code{"cc"} 告诉编译器 condition codes 被修改。

\subsubsection{错误示例：忘记 \code{"memory"} clobber}

考虑一个人为的 store：

\begin{lstlisting}[language=C++]
void broken_store(std::int64_t* p, std::int64_t value)
{
    asm volatile("mov %1, (%0)" : : "r"(p), "r"(value));
}
\end{lstlisting}

汇编确实写了 \code{*p}，但 C++ 优化器不知道这个 asm 会改变任意内存。若周围代码读取或写入同一内存，优化器可能重排。

更清晰的写法是把内存作为输出操作数：

\begin{lstlisting}[language=C++]
void better_store(std::int64_t* p, std::int64_t value)
{
    asm volatile("mov %1, %0"
        : "=m"(*p)
        : "r"(value)
        : "memory");
}
\end{lstlisting}

这里 \code{"=m"(*p)} 明确描述被写的内存对象，\code{"memory"} 进一步作为编译器级内存屏障，防止优化器把其他内存访问跨过这个 asm 移动。

\begin{warningbox}
\code{"memory"} clobber 是编译器屏障，不是 CPU 内存栅栏。它限制编译器重排，但不会自动生成 \instruction{mfence}、\instruction{lfence} 或 locked instruction。多线程内存序问题必须用 C++ atomics 或明确的硬件同步指令解决。
\end{warningbox}

\subsubsection{\code{volatile} 不是万能药}

\code{asm volatile} 的含义接近“这个 asm 有副作用，不要因为输出未使用就删除它，也不要随便合并”。但它不自动说明：

\begin{itemize}
  \item asm 读写了哪些 C++ 变量。
  \item asm 修改了哪些寄存器。
  \item asm 修改了 flags。
  \item asm 访问了内存。
  \item asm 可以作为 CPU fence。
\end{itemize}

所以 inline asm 的正确性来自 operands 和 clobbers，而不是只靠 \code{volatile}。

\subsubsection{常见约束}

\begin{center}
\begin{tabular}{p{0.16\linewidth}p{0.68\linewidth}}
\toprule
约束 & 含义 \\
\midrule
\code{"r"} & 任意通用寄存器。 \\
\code{"m"} & 内存操作数。 \\
\code{"i"} & 编译期立即数。 \\
\code{"a"} & 使用 \register{rax}/\register{eax}/\register{al} 类寄存器。 \\
\code{"=r"} & 只写输出寄存器。 \\
\code{"+r"} & 读写同一个寄存器操作数。 \\
\code{"=&r"} & early-clobber 输出，输出在所有输入读完前就可能被写。 \\
\code{"0"} & 与第 0 个输出操作数使用同一位置。 \\
\code{"cc"} & asm 修改了 flags。 \\
\code{"memory"} & asm 可能读写未显式列出的内存，限制编译器重排。 \\
\bottomrule
\end{tabular}
\end{center}

\begin{deepdive}
inline asm 的约束是你和编译器寄存器分配器之间的合同。你不应该在模板里随便写固定寄存器名，然后忘记告诉编译器。固定寄存器越多，寄存器分配越难，优化器越受限制，错误空间也越大。
\end{deepdive}

\subsection{intrinsics、inline asm 与 \code{.S} 的选择}

现代高性能代码通常优先选择下面顺序：

\begin{center}
\begin{tabular}{p{0.20\linewidth}p{0.34\linewidth}p{0.30\linewidth}}
\toprule
方式 & 优点 & 风险 \\
\midrule
普通 C++ & 可移植、优化器可理解、测试容易。 & 不一定生成你想要的指令。 \\
compiler builtin & 表达特殊语义，如 prefetch、expect、overflow。 & 依赖编译器扩展。 \\
intrinsics & 直接表达 SIMD 指令，寄存器分配仍由编译器做。 & ISA 绑定强，代码可读性下降。 \\
\code{.S} 文件 & 可完全控制函数级汇编和 ABI。 & 维护成本高，调试和 unwind 要额外注意。 \\
inline asm & 可在表达式附近插入特殊指令。 & 最容易和优化器合同写错。 \\
\bottomrule
\end{tabular}
\end{center}

对于 AVX2/AVX-512/FMA/VNNI 这类 AI 算子优化，本书后续会大量使用 intrinsics。原因是 intrinsics 让编译器仍然负责寄存器分配、调用约定、栈框架和调度的一部分，同时你可以明确控制核心 SIMD 操作。

单独 \code{.S} 文件适合：

\begin{itemize}
  \item 你需要完全固定函数入口、寄存器分配和指令布局。
  \item 你正在写非常小的热路径 kernel。
  \item 你需要验证某个 ABI 或机器码现象。
  \item 你愿意维护多 ISA 版本和测试矩阵。
\end{itemize}

inline asm 适合：

\begin{itemize}
  \item 少数没有合适 intrinsic 的特殊指令。
  \item 需要非常局部的硬件交互。
  \item 你已经能准确写出 operands、constraints 和 clobbers。
\end{itemize}

\subsection{调试和审计工作流}

写互操作代码时，永远保留可审计证据。

\subsubsection{编译参数}

\begin{lstlisting}[language=bash]
clang++ -std=c++20 -O2 -g -fno-omit-frame-pointer \
    main.cpp kernels.S -o app
\end{lstlisting}

训练阶段建议：

\begin{itemize}
  \item 用 \code{-g} 保留调试信息。
  \item 初期加 \code{-fno-omit-frame-pointer}，让栈回溯更直观。
  \item 分别看 \code{-O0} 和 \code{-O2/-O3}，理解优化差异。
  \item 对汇编文件保留注释和清晰符号名。
\end{itemize}

\subsubsection{符号审计}

\begin{lstlisting}[language=bash]
nm -C app | rg 'asm_|cpp_'
readelf -Ws app | rg 'asm_|cpp_'
\end{lstlisting}

要检查：

\begin{itemize}
  \item 汇编函数是否真的导出。
  \item C++ 是否引用了预期符号名。
  \item 是否出现未预期的 mangled name。
  \item 符号类型是否是函数。
\end{itemize}

\subsubsection{反汇编审计}

\begin{lstlisting}[language=bash]
objdump -drwC -Mintel app > app.objdump.txt
rg -n '<asm_sum8>|<asm_square_plus>|call.*cpp_square' app.objdump.txt
\end{lstlisting}

要检查：

\begin{itemize}
  \item 参数寄存器使用是否和 ABI 一致。
  \item 栈参数偏移是否正确。
  \item callee-saved 寄存器是否成对保存恢复。
  \item 调用前 \register{rsp} 是否对齐。
  \item 是否出现你没预料到的 PLT 调用或间接跳转。
\end{itemize}

\subsubsection{GDB 单步}

\begin{lstlisting}[language=bash]
gdb ./app
(gdb) break asm_square_plus
(gdb) run
(gdb) disassemble /r asm_square_plus
(gdb) info registers rdi rax rbx rsp rip
(gdb) x/6gx $rsp
(gdb) stepi
(gdb) info registers rdi rax rbx rsp rip
\end{lstlisting}

单步时不要只看源代码行。要看：

\begin{itemize}
  \item \register{rip} 是否进入预期符号。
  \item \register{rsp} 每次 \instruction{push}、\instruction{pop}、\instruction{call}、\instruction{ret} 如何变化。
  \item 返回地址是否位于预期栈槽。
  \item 返回值是否写入 \register{rax} 或 \register{xmm0}。
\end{itemize}

\subsection{sanitizer 和手写汇编的边界}

AddressSanitizer、UndefinedBehaviorSanitizer、ThreadSanitizer 对 C++ 很有价值，但它们不能自动理解你所有手写汇编内存访问。比如汇编里：

\begin{lstlisting}
mov rax, qword ptr [rdi + rcx*8]
\end{lstlisting}

如果 \register{rdi} 指向越界地址，ASan 未必能像检查 C++ 数组访问那样插桩发现。原因是插桩发生在编译器能理解的 IR 层；手写汇编已经绕过了这部分语义。

因此汇编边界测试要更保守：

\begin{itemize}
  \item C++ wrapper 先检查尺寸、空指针和对齐前置条件。
  \item 汇编 kernel 文档明确要求输入合法。
  \item 随机测试使用保护页或额外 padding 检查越界。
  \item 对 debug 版本提供纯 C++ reference path。
  \item 必要时在汇编函数外层做 shadow check，而不是假设 sanitizer 能看见内部。
\end{itemize}

\begin{keyidea}
优化 kernel 的工程边界通常是“外层 C++ 负责合法性和调度，内层汇编负责极少分支的热路径”。不要把复杂错误处理塞进手写汇编热循环。
\end{keyidea}

\subsection{一个接近算子 kernel 的例子：整数点积}

先写标量整数点积，为后续 SIMD 章节做准备。

C++ 声明：

\begin{lstlisting}[language=C++]
extern "C" std::int64_t asm_dot_i32(const std::int32_t* a,
                                    const std::int32_t* b,
                                    std::int64_t n);
\end{lstlisting}

语义：

\begin{lstlisting}
sum = 0
for i in [0, n):
    sum += int64(a[i]) * int64(b[i])
return sum
\end{lstlisting}

ABI 入口：

\begin{lstlisting}
rdi = a
rsi = b
rdx = n
rax = return value
\end{lstlisting}

汇编：

\begin{lstlisting}
.intel_syntax noprefix
.text

.globl asm_dot_i32
.type asm_dot_i32, @function
asm_dot_i32:
    xor rax, rax
    test rdx, rdx
    jle .Ldone

    xor rcx, rcx
.Lloop:
    movsxd r8, dword ptr [rdi + rcx*4]
    movsxd r9, dword ptr [rsi + rcx*4]
    imul r8, r9
    add rax, r8
    inc rcx
    cmp rcx, rdx
    jl .Lloop

.Ldone:
    ret
.size asm_dot_i32, .-asm_dot_i32

.section .note.GNU-stack,"",@progbits
\end{lstlisting}

地址公式：

\begin{lstlisting}
address of a[i] = rdi + rcx * 4
address of b[i] = rsi + rcx * 4
\end{lstlisting}

为什么 scale 是 4？因为 \code{std::int32_t} 元素大小是 4 字节。这个 scale 是 x86 地址模式支持的合法 scale。若元素结构大小是 12 字节，就不能写 \code{[base + index*12]}，必须用 \instruction{lea} 或其他指令拆解。

这段代码还不是高性能点积。它每次循环只做一个元素，乘法和加载串行依赖较多，没有展开，没有 SIMD，也没有处理内存带宽。但它是非常好的互操作训练材料，因为它包含：

\begin{itemize}
  \item 指针参数。
  \item signed load extension。
  \item 循环控制流。
  \item 数组地址公式。
  \item 64 位累加返回值。
\end{itemize}

\subsection{汇编边界的文档模板}

每个手写汇编函数都应有一段接口注释。示例：

\begin{lstlisting}
# std::int64_t asm_dot_i32(const std::int32_t* a,
#                          const std::int32_t* b,
#                          std::int64_t n)
#
# ABI:
#   rdi = a
#   rsi = b
#   rdx = n
#   rax = return sum
#
# Preconditions:
#   a and b point to at least n int32 elements
#   n >= 0
#   pointers may be unaligned
#
# Clobbers:
#   rax, rcx, r8, r9, flags
#
# Callee-saved:
#   none modified
\end{lstlisting}

这不是形式主义。几个月后你优化这个 kernel 时，注释会告诉你哪些寄存器本来可以自由使用，哪些语义由外层保证，哪些行为必须保持。

\subsection{实验 14.1：从零搭建 C++ 调汇编工程}

\begin{labbox}
目标：建立一个最小但规范的 \code{.S + C++ + CMake} 工程，并用工具证明符号、调用和返回值正确。

创建目录：

\begin{lstlisting}[language=bash]
mkdir -p labs/ch14_interop/lab14_1_minimal
cd labs/ch14_interop/lab14_1_minimal
\end{lstlisting}

文件：

\begin{lstlisting}
lab14_1_minimal/
    CMakeLists.txt
    main.cpp
    asm_arith.S
\end{lstlisting}

任务：

\begin{enumerate}
  \item 在 C++ 中声明 3 个函数：
\begin{lstlisting}[language=C++]
extern "C" long asm_add2(long a, long b);
extern "C" long asm_sub2(long a, long b);
extern "C" long asm_mix4(long a, long b, long c, long d);
\end{lstlisting}
  \item 在 \code{asm_arith.S} 中实现它们。
  \item \code{asm_mix4} 的语义为 \code{(a + b) * 3 - (c - d)}，可以使用 \instruction{lea}、\instruction{add}、\instruction{sub}。
  \item 写至少 20 个确定性测试，覆盖正数、负数、零和较大值。
  \item 构建并运行。
  \item 用 \code{nm -C} 证明符号存在。
  \item 用 \code{objdump -drwC -Mintel} 找到 3 个函数的机器码。
\end{enumerate}

推荐命令：

\begin{lstlisting}[language=bash]
cmake -S . -B build -DCMAKE_BUILD_TYPE=RelWithDebInfo
cmake --build build -j
./build/lab14_1_minimal
nm -C build/lab14_1_minimal | rg 'asm_'
objdump -drwC -Mintel build/lab14_1_minimal > lab14_1.objdump.txt
\end{lstlisting}

交付物：

\begin{enumerate}
  \item 源码。
  \item 运行输出。
  \item \code{nm} 输出中 3 个汇编符号的行。
  \item 每个函数的反汇编。
  \item 一张表：源代码参数、ABI 寄存器、汇编使用位置。
\end{enumerate}
\end{labbox}

\subsection{实验 14.2：符号、mangling 和 relocation 审计}

\begin{labbox}
目标：理解 C++ 符号名、\code{extern "C"} 和链接器重定位。

创建两个版本：

\begin{lstlisting}
with_c_linkage.cpp
without_c_linkage.cpp
caller.cpp
\end{lstlisting}

\code{with_c_linkage.cpp}：

\begin{lstlisting}[language=C++]
extern "C" int add_c(int a, int b)
{
    return a + b;
}
\end{lstlisting}

\code{without_c_linkage.cpp}：

\begin{lstlisting}[language=C++]
int add_cpp(int a, int b)
{
    return a + b;
}
\end{lstlisting}

\code{caller.cpp}：

\begin{lstlisting}[language=C++]
extern "C" int add_c(int, int);
int add_cpp(int, int);

extern "C" int call_both(int x)
{
    return add_c(x, 1) + add_cpp(x, 2);
}
\end{lstlisting}

任务：

\begin{enumerate}
  \item 分别编译成 \code{.o}。
  \item 用 \code{nm} 和 \code{nm -C} 比较符号名。
  \item 用 \code{objdump -drwC -Mintel caller.o} 查看未链接目标文件。
  \item 用 \code{readelf -r caller.o} 查看 relocation。
  \item 链接成可执行文件后，再看最终 \instruction{call} 的相对位移。
\end{enumerate}

命令：

\begin{lstlisting}[language=bash]
clang++ -std=c++20 -O2 -c with_c_linkage.cpp -o with_c_linkage.o
clang++ -std=c++20 -O2 -c without_c_linkage.cpp -o without_c_linkage.o
clang++ -std=c++20 -O2 -c caller.cpp -o caller.o
nm caller.o
nm -C caller.o
objdump -drwC -Mintel caller.o > caller.objdump.txt
readelf -r caller.o > caller.reloc.txt
\end{lstlisting}

交付物：

\begin{enumerate}
  \item \code{add_c} 和 \code{add_cpp} 的原始符号名。
  \item \code{call_both} 中两个 \instruction{call} 的 relocation 记录。
  \item 解释 \code{R_X86_64_PLT32} 或你机器上出现的 relocation 类型。
  \item 链接后任选一个 \instruction{call}，用 \code{target - next_rip} 计算其 \code{rel32}。
  \item 一段 500 字总结：为什么目标文件反汇编不能只看机器码字节。
\end{enumerate}
\end{labbox}

\subsection{实验 14.3：ABI 安全的 8 参数汇编函数}

\begin{labbox}
目标：实现并严格测试一个包含栈参数的汇编函数。

函数签名：

\begin{lstlisting}[language=C++]
extern "C" long asm_weighted_sum8(long a0, long a1, long a2, long a3,
                                  long a4, long a5, long a6, long a7);
\end{lstlisting}

语义：

\begin{lstlisting}
return a0*1 + a1*2 + a2*3 + a3*4
     + a4*5 + a5*6 + a6*7 + a7*8
\end{lstlisting}

任务：

\begin{enumerate}
  \item 写 ABI 参数位置表。
  \item 实现汇编函数。第 7、8 个参数必须从栈读取。
  \item 写 C++ reference 函数。
  \item 写 100000 轮随机测试，限制输入范围避免 signed overflow。
  \item 用 GDB 在函数入口断点，记录 \register{rsp}、\code{[rsp]}、\code{[rsp + 8]}、\code{[rsp + 16]}。
  \item 用 \code{objdump} 标注每个参数在哪里被使用。
\end{enumerate}

提示：可以用 \instruction{lea} 组合小常数乘法，但要保证结果语义清楚。例如 \code{x*3} 可以写成：

\begin{lstlisting}
lea rax, [rdi + rdi*2]
\end{lstlisting}

交付物：

\begin{enumerate}
  \item 参数映射表。
  \item 完整汇编。
  \item 测试源码和运行结果。
  \item GDB 栈截图或文字记录。
  \item 解释为什么 \code{a6} 是 \code{[rsp + 8]} 而不是 \code{[rsp]}。
\end{enumerate}
\end{labbox}

\subsection{实验 14.4：汇编调用 C++ helper}

\begin{labbox}
目标：写一个 non-leaf 汇编函数，正确处理 callee-saved 寄存器和栈对齐。

C++ helper：

\begin{lstlisting}[language=C++]
extern "C" long cpp_mul_add(long x, long y, long z)
{
    return x * y + z;
}
\end{lstlisting}

汇编函数签名：

\begin{lstlisting}[language=C++]
extern "C" long asm_call_helper(long x, long y, long z, long bias);
\end{lstlisting}

语义：

\begin{lstlisting}
return cpp_mul_add(x, y, z) + bias
\end{lstlisting}

任务：

\begin{enumerate}
  \item 先故意写一个错误版本：把 \code{bias} 放进 \register{rbx}，调用 helper 后不恢复。
  \item 构造 C++ 调用者，让调用者在调用前后有一个长期存活变量，观察错误是否暴露。
  \item 修复 \register{rbx} 保存恢复。
  \item 检查调用 \code{cpp_mul_add} 前 \register{rsp} 是否 16 字节对齐。
  \item 用 GDB 单步过 \instruction{call}，记录 \register{rsp} 和返回地址。
\end{enumerate}

交付物：

\begin{enumerate}
  \item 错误版本和修复版本的汇编。
  \item 为什么错误版本违反 ABI。
  \item 栈对齐计算过程。
  \item GDB 记录。
  \item 测试结果。
\end{enumerate}
\end{labbox}

\subsection{实验 14.5：inline asm 约束修复}

\begin{labbox}
目标：理解 inline asm 不是文本替换，而是优化器合同。

给定错误代码：

\begin{lstlisting}[language=C++]
#include <cstdint>

std::int64_t broken_add(std::int64_t a, std::int64_t b)
{
    std::int64_t result = a;
    asm("add %1, %0" : : "r"(result), "r"(b));
    return result;
}

void broken_store(std::int64_t* p, std::int64_t value)
{
    asm volatile("mov %1, (%0)" : : "r"(p), "r"(value));
}
\end{lstlisting}

任务：

\begin{enumerate}
  \item 在 \code{-O0} 和 \code{-O3} 下分别编译运行，记录现象。
  \item 查看反汇编，解释为什么错误版本可能“偶尔看起来能跑”。
  \item 修复 \code{broken_add}，正确使用读写输出和 \code{"cc"} clobber。
  \item 修复 \code{broken_store}，正确描述内存写和 \code{"memory"} clobber。
  \item 写测试证明修复版本在 \code{-O3} 下稳定正确。
\end{enumerate}

交付物：

\begin{enumerate}
  \item 错误版本和修复版本源码。
  \item \code{-O0}、\code{-O3} 反汇编对比。
  \item 每个 constraint 的解释。
  \item 说明 \code{volatile}、\code{"memory"}、\code{"cc"} 各自解决什么问题，不能解决什么问题。
\end{enumerate}
\end{labbox}

\subsection{实验 14.6：标量点积 kernel 的互操作审计}

\begin{labbox}
目标：实现一个接近算子 kernel 形态的标量点积，并建立 correctness report。

函数：

\begin{lstlisting}[language=C++]
extern "C" std::int64_t asm_dot_i32(const std::int32_t* a,
                                    const std::int32_t* b,
                                    std::int64_t n);
\end{lstlisting}

任务：

\begin{enumerate}
  \item 写 C++ reference 实现。
  \item 写汇编实现。
  \item 测试 \code{n = 0, 1, 2, 15, 16, 17, 1024}。
  \item 随机生成数组，做 10000 轮测试。
  \item 用 \code{objdump} 标注循环基本块、回边、数组地址公式。
  \item 用 \code{perf stat} 初步测量运行时间、instructions、cycles；如果权限不足，记录错误信息并使用 \code{/usr/bin/time}。
\end{enumerate}

交付物：

\begin{enumerate}
  \item 源码和构建脚本。
  \item 正确性测试输出。
  \item 循环反汇编标注。
  \item 地址公式说明。
  \item 初步性能表。
  \item 说明为什么这个版本还不是高性能实现，列出至少 5 个后续优化方向。
\end{enumerate}
\end{labbox}

\subsection{本章作业}

\begin{exercisebox}
\subsubsection*{作业 1：接口契约报告}

选择你在实验中写的任意 3 个汇编函数，为每个函数写一份接口契约报告，必须包含：

\begin{itemize}
  \item C++ 声明。
  \item 链接符号名。
  \item 参数位置表。
  \item 返回值位置。
  \item 修改的 caller-saved 寄存器。
  \item 是否修改 callee-saved 寄存器；如果修改，如何保存恢复。
  \item 是否调用其他函数；如果调用，如何保证栈对齐。
  \item 前置条件和未定义行为边界。
\end{itemize}

\subsubsection*{作业 2：修复破坏 ABI 的汇编}

下面函数被 C++ 调用：

\begin{lstlisting}
.intel_syntax noprefix
.text
.globl broken_kernel
broken_kernel:
    mov r12, rdi
    mov r13, rsi
    call helper
    add rax, r12
    add rax, r13
    ret
\end{lstlisting}

要求：

\begin{enumerate}
  \item 找出所有 ABI 问题。
  \item 说明 \register{r12}、\register{r13} 属于 caller-saved 还是 callee-saved。
  \item 修复保存恢复。
  \item 修复栈对齐。
  \item 给出完整汇编。
  \item 用栈布局图证明你的修复正确。
\end{enumerate}

\subsubsection*{作业 3：重定位计算}

给定未链接目标文件片段：

\begin{lstlisting}
0000000000000000 <foo>:
   0: e8 00 00 00 00    call 5 <foo+0x5>
        1: R_X86_64_PLT32 bar-0x4
   5: c3                ret
\end{lstlisting}

链接后：

\begin{lstlisting}
foo address = 0x401100
bar address = 0x401180
\end{lstlisting}

要求：

\begin{enumerate}
  \item 写出 \instruction{call} 指令地址。
  \item 写出 next \register{rip}。
  \item 计算 \code{rel32}。
  \item 写出小端序 4 字节。
  \item 解释为什么 relocation addend 常见为 \code{-4}。
\end{enumerate}

\subsubsection*{作业 4：inline asm 审计}

阅读下面代码：

\begin{lstlisting}[language=C++]
std::uint64_t rdtsc_bad()
{
    std::uint32_t lo = 0;
    std::uint32_t hi = 0;
    asm volatile("rdtsc");
    return (static_cast<std::uint64_t>(hi) << 32) | lo;
}
\end{lstlisting}

要求：

\begin{enumerate}
  \item 解释为什么它错误。
  \item 写出正确输出约束，让 \register{eax} 和 \register{edx} 分别写入 \code{lo}、\code{hi}。
  \item 说明 \instruction{rdtsc} 的乱序风险。
  \item 查阅本机反汇编，确认约束是否生成预期寄存器。
  \item 说明为什么 benchmark 不能只靠裸 \instruction{rdtsc}。
\end{enumerate}

\subsubsection*{作业 5：互操作小项目}

完成一个小项目：

\begin{lstlisting}
interop_project/
    CMakeLists.txt
    include/kernels.hpp
    src/main.cpp
    src/reference.cpp
    src/kernels.S
    tests/test_kernels.cpp
    report.md
\end{lstlisting}

功能：

\begin{itemize}
  \item \code{asm_sum8}：8 参数加权求和。
  \item \code{asm_make_pair}：返回两个 64 位字段的小结构体。
  \item \code{asm_dot_i32}：标量点积。
  \item \code{asm_call_helper}：汇编调用 C++ helper。
\end{itemize}

质量要求：

\begin{enumerate}
  \item 每个函数都有 reference 实现。
  \item 每个函数都有确定性测试和随机测试。
  \item 所有汇编函数都有接口注释。
  \item 报告中包含 \code{nm}、\code{readelf -r}、\code{objdump} 证据。
  \item 报告中至少分析 2 个栈布局。
  \item 报告中至少解释 1 个 relocation。
  \item 报告中明确列出 sanitizer 能覆盖什么、不能覆盖什么。
\end{enumerate}

\subsubsection*{评分标准}

本章作业按下面标准评价：

\begin{enumerate}
  \item 能否从 C++ 签名正确推导 ABI 参数和返回值位置。
  \item 汇编符号、\code{.type}、\code{.size}、section 是否规范。
  \item 是否正确处理 caller-saved、callee-saved 和栈对齐。
  \item 是否能用 \code{nm}、\code{readelf}、\code{objdump}、GDB 证明结论。
  \item inline asm 是否准确描述输入、输出和 clobbers。
  \item 测试是否覆盖边界值、随机值和跨优化级别行为。
  \item 报告是否区分语言层语义、ABI 层契约、目标文件层机制和 CPU 执行行为。
\end{enumerate}
\end{exercisebox}

\subsection{本章小结}

本章把 System V AMD64 ABI 放进了真实工程：

\begin{lstlisting}
C++ declaration:
    type, language linkage, call site generation

assembly source:
    symbol definition, sections, prologue/epilogue, instructions

ABI:
    parameter locations, return locations, register ownership, stack alignment

object file:
    symbols, sections, relocations

tools:
    nm, readelf, objdump, gdb, perf

quality:
    deterministic tests, randomized tests, binary-interface report
\end{lstlisting}

你现在应该能够从一个 C++ 函数声明出发，写出 ABI 映射表，再写出手写汇编实现，并用工具证明它真的按预期链接和运行。你也应该开始对 inline asm 保持谨慎：它不是简单的“在 C++ 里写汇编”，而是和优化器签订一份非常精确的合同。

后面的章节会进入测量、编译器优化、微架构和 SIMD。本章建立的互操作能力会反复使用：你会写小 kernel、反汇编审计、测量性能、解释瓶颈，再逐步把标量代码推进到 AVX、FMA、VNNI 和 AI 算子优化。
